% Working replacement fragment.  Not integrated or claimed closed.
% The external Gaussian Hecke Siegel--Walfisz input below remains subject
% to the independent source/proof review in published-input-audit/.
\subsection{Gaussian ray weights and their progression main terms}
\label{subsec:gaussian-progression-mainterms}

Let $E=\Q(i)$ and let $H$ be a fixed positive integer divisible by $16$.
Let $K/E$ be the ray class field for $H\Z[i]$.  This field is Galois over
$\Q$, since complex conjugation preserves both the modulus and its
principal ray subgroup.  Put
\[
 A=\operatorname{Gal}(K/E)
   \simeq (\Z[i]/H\Z[i])^\times/\{1,-1,i,-i\},
 \qquad G=\operatorname{Gal}(K/\Q).
\]
We use arithmetic Frobenius throughout.  Let $C\subseteq A$ be invariant
under complex conjugation, so that $C$ is a union of conjugacy classes in
$G$.  Define $\Lambda'_C(p)=\log p$ for unramified rational primes whose
Frobenius lies in $C$, and put $\Lambda'_C(n)=0$ otherwise.  In particular
this function is supported on rational primes which split in $E$.

For ideals coprime to $H$ let $[\mathfrak a]\in A$ denote their ray class,
and set
\[
 T_C(n)=\sum_{\substack{\mathrm N\mathfrak a=n\\(\mathfrak a,H)=1}}
       \Lambda_E(\mathfrak a)1_C([\mathfrak a]).
\]
Here $\Lambda_E(\mathfrak p^j)=\log\mathrm N\mathfrak p$ and it is zero
on all other ideals.  Every unramified split rational prime has two
conjugate prime ideals of norm $p$, so
\begin{equation}\label{eq:gaussian-ray-prime-factor}
 T_C(p)=2\Lambda'_C(p)\qquad(p\nmid2H).
\end{equation}
The remaining difference between $T_C$ and $2\Lambda'_C$ is supported on
rational prime powers of exponent at least two and the finitely many
excluded primes.  Its total absolute value up to $X$ is
$O_H(\sqrt X\log^2(2X))$.  This bound follows by summing the at most two
prime ideals above each rational prime and their powers; no density
theorem is used in this comparison.

For $\chi\in\widehat A$ put
\[
 c_C(\chi)=|A|^{-1}\sum_{a\in C}\overline{\chi(a)},\qquad
 T_\chi(n)=\sum_{\mathrm N\mathfrak a=n}
                  \Lambda_E(\mathfrak a)\chi(\mathfrak a),
\]
where each character is extended by zero on ideals not coprime to $H$.
Finite Fourier inversion gives the exact identity
\begin{equation}\label{eq:gaussian-ray-fourier}
                 T_C(n)=\sum_{\chi\in\widehat A}c_C(\chi)T_\chi(n).
\end{equation}
The field and all these characters are fixed before $X$ tends to infinity.

The progression input needed below is the following fixed-field Hecke
Siegel--Walfisz assertion.  For every $B,D>0$, every fixed finite-order
ray character $\chi$ of $E$, and every Dirichlet character $\rho$ of
modulus $q\leq(\log X)^B$,
\begin{equation}\label{eq:gaussian-hecke-sw-input}
 \sum_{\mathrm N\mathfrak a\leq Y}
       \Lambda_E(\mathfrak a)\chi(\mathfrak a)
                         \rho(\mathrm N\mathfrak a)
  = r_{\chi,\rho}Y+O_{H,B,D}(X(\log X)^{-D})
 \quad(0\leq Y\leq X),
\end{equation}
where $r_{\chi,\rho}=1$ if the induced Hecke character
$\chi(\rho\circ\mathrm N)$ is trivial and is zero otherwise.
The omitted Euler factors only change the sum by a polylogarithmic
quantity.  The constant may be ineffective.  We require uniformity in
$\rho,q,Y$, but not in $H$.
% SOURCE GATE: insert the independently checked published source and
% explicit correction proof before integrating this fragment.

\begin{lemma}[Progression density from the fixed-field input]
\label{lem:gaussian-ray-progression}
Assume \eqref{eq:gaussian-hecke-sw-input}.  For $(a,q)=1$ define
$J_q=K\cap\Q(\zeta_q)$ and let $\sigma_a(\zeta_q)=\zeta_q^a$.
For fixed $B,D>0$, uniformly for $q\leq(\log X)^B$ and $0\leq Y\leq X$,
\[
 \sum_{\substack{n\leq Y\\ n\equiv a\pmod q}}\Lambda'_C(n)
 =\frac{\#\{g\in C:g|_{J_q}=\sigma_a|_{J_q}\}}
             {[K\Q(\zeta_q):\Q]}Y
       +O_{H,B,D}(X(\log X)^{-D}).
\]
\end{lemma}
\begin{proof}
Orthogonality of Dirichlet characters and
\eqref{eq:gaussian-ray-fourier} show that the coefficient of $Y$ for
$T_C$ is
\[
 \frac1{\varphi(q)}\sum_{\chi\in\widehat A}c_C(\chi)
       \sum_{\rho\bmod q}\overline{\rho(a)}r_{\chi,\rho}.
\]
The finite sum over $\chi$ has size depending only on $H$, while the
average over $\rho$ cancels its own factor $\varphi(q)$ in the error.

To evaluate this coefficient without a convention about inverse Artin
symbols, put $L=K\Q(\zeta_q)$ and $B_q=\operatorname{Gal}(L/E)$.
Restriction identifies $B_q$ with
\[
 \{(g,s)\in A\times\operatorname{Gal}(\Q(\zeta_q)/\Q):
                       g|_{J_q}=s|_{J_q}\}.
\]
The characters $\chi(g)\rho(s)$ restrict to characters of $B_q$.
Their average on $B_q$ is one precisely when their pullback by Artin
reciprocity is the trivial Hecke character; otherwise it is zero.
This is exactly $r_{\chi,\rho}$.  Expanding $1_C(g)$ and
$1_{s=\sigma_a}$ by Fourier inversion therefore identifies the displayed
coefficient with the proportion of elements of $B_q$ satisfying
$g\in C$ and $s=\sigma_a$.  Its numerator is the numerator in the
lemma and its denominator is $|B_q|=[L:E]$.
Dividing by two using \eqref{eq:gaussian-ray-prime-factor} changes the
denominator to $[L:\Q]$.  The prime-power error is smaller than the stated
error.  This proves the formula.
\end{proof}

Choose a fixed integer $W$ divisible by the conductor of $K^{\mathrm{ab}}$.
It is in particular divisible by $4$.
For a unit $b\pmod W$ put
\[
 \mu_C(b)=
 \frac{\#\{g\in C:g|_{K^{\mathrm{ab}}}=\sigma_b|_{K^{\mathrm{ab}}}\}}
      {[K:K^{\mathrm{ab}}]},\quad
 f_{C,b,W}(n)=\frac{\varphi(W)}W\Lambda'_C(Wn+b),\quad
 \lambda_{b,W}(n)=\frac{\varphi(W)}W\Lambda'(Wn+b).
\]
The prime-supported ordinary von Mangoldt function is denoted by
$\Lambda'$.  Replacing it by the usual $\Lambda$ costs only the same
negligible prime-power error.
The centered weight is
\begin{equation}\label{eq:gaussian-centered-weight}
       h_{C,b,W}=f_{C,b,W}-\mu_C(b)\lambda_{b,W}.
\end{equation}
Here $0\leq\mu_C(b)\leq1$, so
$|h_{C,b,W}(n)|\leq\lambda_{b,W}(n)=O_W(\log(2N))$ on $[N]$.

\begin{corollary}[Centered sums on progressions]
\label{cor:gaussian-centered-progressions}
Assume \eqref{eq:gaussian-hecke-sw-input}.  For every $D>0$,
\[
       \sum_{n\in P}h_{C,b,W}(n)=O_{H,W,D}(N(\log N)^{-D})
\]
uniformly over arithmetic progressions $P$ contained in $[N]$.
\end{corollary}
\begin{proof}
Write $P$ as a residue class modulo $q$ intersected with an interval.
If $q>(\log N)^{D+3}$, its cardinality is at most
$N/(\log N)^{D+3}+1$, and the pointwise logarithmic bound proves the
claim directly.  Otherwise the condition $n\in P$ becomes one residue
class $a\pmod{Wq}$ and an interval for $Wn+b$.

If $(a,Wq)>1$, every prime in this residue class divides $Wq$.  There
are at most $O(\log(Wq))$ such primes, with total logarithmic weight
polylogarithmic in $N$, which is negligible here.  Suppose instead that
$(a,Wq)=1$.
Since $\Q(\zeta_{Wq})$ is abelian and already contains
$K^{\mathrm{ab}}$, one has
$K\cap\Q(\zeta_{Wq})=K^{\mathrm{ab}}$.
Also $a\equiv b\pmod W$, so the lemma's main coefficient is
\[
   \frac{\mu_C(b)}{\varphi(Wq)}.
\]
For the ordinary prime weight, the Siegel--Walfisz main coefficient
is $1/\varphi(Wq)$.  The two terms of
\eqref{eq:gaussian-centered-weight} therefore have equal main terms.
Apply the progression lemma at both interval endpoints, with an error
exponent larger than $D$ and with $X=O_W(N)$, and subtract.  This proves
the claimed bound.  All moduli used are at most a fixed power of
$\log N$, and $W$ is fixed throughout this argument.
\end{proof}
